%=====================================================================
\section{Learning the Procedure from a Handful of Gold Annotations}
\label{sec:learning}

All learning here changes $\Pi$ (text, code, thresholds, registry) and never a model. For \umr documents $D$ we measure
\begin{equation}
J(\Pi; D) = \frac{1}{|D|}\sum_{d \in D} \tfrac{1}{2}\bigl(F_d + C_d\bigr),
\label{eq:objective}
\end{equation}
where $F_d$ is the sentence-graph F-score and $C_d$ the comprehensive score that also covers the document level \citep{sun2024ancast}; in batch mode the score is the task's main metric on a set of sentences. Errors are diagnosed on \emph{training} data; \emph{regression} data (earlier training data) must not get worse; \emph{development} data are measured but never shown for diagnosis, though acceptance uses their scores, so they are model-selection data. Test data are excluded from edit generation and acceptance.

\paragraph{Threshold tuning.}
\label{sec:calibration}
Rules read probabilities against thresholds set for one decision model; another model spreads its probability mass differently, so a rule fires too often or too rarely. What matters is the decision a rule takes \citep[cf.][]{zhao2021calibrate}, so we fit $\theta$ to $J$ by coordinate descent, sweeping one threshold at a time over a few values and keeping a value that raises $J$ by more than 0.05.

\paragraph{Routing.}
\label{sec:routing}
A closed answer has a confidence $c=\max_o p(o)$. Questions with $c<\tau$ are asked again, with the same state and options, to a stronger answerer, an LLM forced to pick an option or, in a cold start, a person \exper{E5}; nothing else changes.

\paragraph{Reflection.}
\label{sec:reflect}
A round on a training document runs four steps with a strong LLM as the \emph{reflector}. \textbf{Find}: from the gold and predicted annotations, the alignment, the judge's failed units and the trace, list each difference, the responsible node and the kind of cause. \textbf{Diagnose}: propose ranked changes per place in the procedure. \textbf{Edit}: write anchored edits. \textbf{Accept}: apply them to a copy of $\Pi$, the next frozen version. In fragment mode a parse takes minutes and dollars, so acceptance is \emph{unmeasured}: an edit is accepted if its anchor matches once and the skill still loads, and the next document tests the round. In closed mode a parse takes seconds, so acceptance is \emph{measured}: proposals are applied one at a time and the training, regression and development documents re-parsed after each. An edit is kept if $J$ rises by at least 0.1, mean $F$ and $C$ fall by at most 0.3, no document falls by more than $\max(0.8, 2\Delta J)$ (capped at 2.5), and the regression documents fall by at most 0.3; a rejected edit is sent back once with its measurement.

\paragraph{Measured reflection in small batches.}
A document per round gives the reflector hundreds of problems at once and few rounds before the gold is used up. For the new runs we reflect on five gold sentences per round (Algorithm~\ref{alg:reflect}, appendix). The reflector reads, for each sentence, the gold and predicted annotations, the scorer's findings and every question asked with its top options and probabilities, together with the editable files and the edits tried before, and proposes at most three edits, each an exact find-and-replace in one file. Each edit is measured on the batch, a regression set (the two previous batches) and a fixed development set, and kept if the mean score rises by at least 0.3, the development score falls by at most 0.5 and the regression score by at most 1.0. Rejected edits go back once with their measured effect. Decisions are cached, so a trial pays only for the questions an edit changes.

\paragraph{Guards.} The reflector never sees test data, and test scores come only from frozen versions. Batch-mode edits are checked for text copied from the test sentences (six Chinese characters or four English words in a row; none found). The closed document-mode check covers six-character Chinese substrings of a configured corpus subset, an indirect use of test text whose coverage remains to be audited \exper{E13}; in fragment mode it was a convention \todo{run the check over the fragment-mode versions}.
